\documentclass[11pt]{article}
\usepackage[T1]{fontenc}
\usepackage[utf8]{inputenc}
\usepackage{lmodern}
\usepackage[margin=1in]{geometry}
\usepackage{amsmath,amssymb,amsthm,mathtools}
\usepackage{enumitem,booktabs,tabularx,array}
\usepackage{microtype}
\usepackage[hidelinks]{hyperref}
\usepackage{tikz-cd}
\newtheorem{theorem}{Theorem}[section]
\newtheorem{proposition}[theorem]{Proposition}
\newtheorem{definition}[theorem]{Definition}
\theoremstyle{remark}
\newtheorem{remark}[theorem]{Remark}
\newcommand{\RR}{\mathbb R}
\newcommand{\QQ}{\mathbb Q}
\newcommand{\ZZ}{\mathbb Z}
\newcommand{\NN}{\mathbb N}
\newcommand{\Sym}{\operatorname{Sym}}
\newcommand{\rank}{\operatorname{rank}}
\newcommand{\tr}{\operatorname{tr}}
\newcommand{\spanlin}{\operatorname{span}}
\newcommand{\diag}{\operatorname{diag}}
\newcommand{\one}{\mathbf 1}
\newcommand{\PSD}{\mathbb S_+}
\newcommand{\OO}{\mathcal O}
\newcommand{\QQF}{\mathcal Q_{\mathrm F}}
\newcommand{\QQG}{\mathcal Q_{\mathrm G}}
\setlist{nosep,leftmargin=*}
\setlength{\parskip}{3pt}
\title{Gate A assessment of two phenomenology candidates}
\author{Technical assessment and corrected proof seeds}
\date{3 November 2026}
\begin{document}
\maketitle
\begin{abstract}
This report tests two concrete candidates against the first research gate in the accepted project direction. The comparative candidate relates a finite positive realization to finite-prefix predictive-state growth. The learning candidate reconstructs a response state from one scalar sensor on a sparsely controlled two-factor linear learner, and compares exact gradient flow with ordinary gradient descent. Both yield precise, noncircular mathematics under strong, visible access and model assumptions. Neither currently establishes the genuine source-level residual required for a new central Gate A result.

The preferred first proof seed, if work continues, is the access-sensitive learning theorem: a connected bipartite graph with one observed edge yields a Gram quotient modulo one signed diagonal under flow, but the complete Gram under gradient descent when there are at least three vertices. We give the proof, a sharp mixed-jet bound, finite endpoint experiments, the two-vertex exception, and native-width consequences. These are a narrow application of classical observability, graph combinatorics, and Gram geometry, not a new general realization principle or a theorem of representation formation. The comparative growth theorem is retained as a secondary control. Its finiteness question is subsumed by stronger weighted-automata results.

The disposition is therefore \textbf{Gate A not yet passed as a new research contribution}. Mathematical viability and novelty are separate. This assessment does not amend the accepted contract, select learning as the definition of phenomenology, or discharge the programme's broader comparative, generative, native, empirical, and cross-instance obligations. This deliverable is intentionally uncompiled \LaTeX{} source.
\end{abstract}
\tableofcontents

\section{Decision and scope}
\label{sec:decision}

Gate A asks for a tractable theorem-sized question, an independently interpretable experimental presentation, a concrete premise/conclusion obligation not already obtained from the strongest predecessor, and explicit controls. Stating a quotient, a compatible fibre, or a generic factorization is insufficient. A correct theorem can still fail this research gate when its claimed contribution is an elementary specialization of known work.

\paragraph{Disposition of the two candidates.}
\begin{enumerate}
\item \textbf{Finite positive realization versus predictive growth.} The mechanisms really are unlike: a finite hidden stochastic implementation and the canonical normalized future-law predictor have different carriers, costs, and universal roles. The repaired candidate is mathematically useful. Its finite-orbit decision is covered by stronger determinization theory; the additional polynomial degree is an elementary virtually abelian orbit calculation whose novelty is not established. Retain it as a control or secondary result, not the missing central comparative theorem.
\item \textbf{Sparse learning responses under flow and gradient descent.} This has the stronger operationally explicit proof seed: known graph connectedness and stated interventions yield an exact observable-space calculation, finite executable probe words, a sharp flow-jet horizon, an information comparison, and a concrete failure of update descent. But it compares two materially different learning regimes \emph{inside the same Gram/observable-module reconstruction method}. It is not a comparison between FCA, predictive filtering, and native factorization. It is eligible for the contract's narrower learning branch, which expressly permits a regime comparison, but does not by itself discharge the full comparative Gate C. The source-level residual remains unestablished.
\end{enumerate}

\paragraph{First attempt if one is pursued.}
Use Theorem~\ref{thm:learning} as the exact statement to challenge first, including its finite endpoint realization in Proposition~\ref{prop:endpoints}. The unresolved research question is whether this sparse-sensor, graph-restricted, optimizer-sensitive classification and its quantitative boundary contribute more than a short corollary of the strongest control and matrix-learning sources. Do not replace that question by a claim that the general observable-algebra construction is new. If an existing theorem or a routine specialization already supplies the entire conclusion, classify this proof as a benchmark and reformulate the research problem before expanding the framework.

The central question remains how observations generate, expose, or select useful structure, and why genuinely different constructions agree or compare with loss. Logical incidence, temporal prediction, positive realization, learning, and later physical or computational instances remain distinct tests. Neither failure here is a no-go theorem for categorical phenomenology. Neither success here would establish a universal theory of emergence.

\section{Baseline and comparison languages}
\label{sec:baseline}

The two-sided project already has substantial results. They are dependencies, not available novelty.
\begin{itemize}
\item A supplied algebraic doctrine gives a free extension of response profiles and a generated image, with the stated initial/terminal properties. The Boolean row quotient versus meet/FCA completion, preservation of generated relations, and test-change obstructions are already in the recovered ledger. Incidence/FCA has its own established concept-generation doctrine; neither candidate proves a new theorem for that doctrine \cite{FCA}.
\item Predictive reconstruction already separates filtering, reachable restriction, and observable minimization. Finite hidden generators can have infinite predictors. Hiding recorded outcomes need not induce a map on old predictive classes; filtering must be performed again.
\item Finite-rank Hankel reconstruction, generated and test-positive cones, and invariant polyhedral or simplicial cone criteria are established dependencies. Linear sufficiency is different from positive or native-resource sufficiency.
\item The completed cylinder and cross-doctrine results already give actual normalization spans, coherent transport under their stated conditions, and local exactness criteria. Actual joint preparations, tests, comparison directions, and mixed-pullback conditions cannot be replaced by a convenient formal envelope.
\item The recent native reports already distinguish public behavior from retained mechanisms, fixed-tensor assembly from independent sequential assembly, and feasibility from optimization and gradient transport. Neither candidate below proves those harder native milestones.
\end{itemize}
The recovered ledger is a source-recovery record, not a new formal verification of all earlier exploratory claims \cite{Ledger,Contract}.

\subsection{Which language earns its use}

For the comparative candidate, response pairings and ordered linear realization keep preparations, continuations, positivity, and hidden-state cost visible. Weighted automata give the strongest existing algorithmic account of projective residuals. Group theory computes the additional growth degree. A two-sided presentation is useful bookkeeping, but no new Isbell completion is required to prove the result.

For the learning candidate, a controlled finite-dimensional matrix module is the shortest faithful proof language. It represents the permitted sequential interventions and the scalar sensor directly. Parametric maps and lenses describe the architecture, metric, reverse derivatives, loss, and optimizer; they do not compute the observational quotient automatically. A relative operational presentation could later organize composition and replacement, but supplying it here would not strengthen the theorem unless a native consequence or comparison with the existing cylinder spans is actually proved.

Stacey's comparison of concrete structures, input/output tests, and forcing is a methodological precedent rather than a source of either numerical formula. Garner's free cylinder completion and lax/colax preservation theorems already cover their stated formal role. CALF's closedness and consistency conditions already control the induced algebra on its wrapper image; generic reconstruction or learning-by-factorization cannot be claimed again \cite{Stacey,Garner,CALF}. No common-language superiority or necessity theorem is established by this report.

\subsection{Parametric and relative controls}
\label{sec:relative}

The operational gradient-learning literature includes stateful optimizers, not only a position-paper architecture language \cite{GradientLearning,Lenses}. In the forward right-action convention, a parameter cell from $(P,f)$ to $(P',f')$ has $\alpha:P\to P'$ and $f=f'(1\odot\alpha)$. The opposite local orientation used by CDL must be translated before comparing presentations. Its Theorem G.10 and Conjecture G.12 have different status; a conjectural extension equivalence is not an established observation bridge \cite{CDL,Contract}.

The accepted relative-computation dependencies retain their full scope:
\begin{description}[style=nextline,leftmargin=0pt]
\item[O1] Under its finite augmented-rig and finitary cartesian-polynomial occurrence hypotheses, relative displays, lift preservation, and Beck--Chevalley conditions, O1 constructs the crossing and coherent modifications. It is stronger than assuming a distributive law. Its necessity claim has its stated occurrence and realizable-projection bounds.
\item[O5] Local certificates produce finite native coupling spans and contextual replacement for each admitted finite typed client, with declared observation and resource comparisons.
\item[O7] Joint law alone can lose zero-probability possibilities. The repaired joint-law-and-possibility semantics has a covered-refinement characterization by native spans and a fixed-endpoint local reflection with coherent lax comparison. It is not executable recovery of all hidden witnesses or a global equivalence.
\end{description}
These are source-level dependencies of the accepted contract, not freshly reproved here \cite{Relative,Contract}. Neither candidate supplies the O1 occurrence presentation, the O5 coupling data, the O7 possibility witnesses, or a pure-compatible multiplicative comparison to an observational relative pseudopromonad. Consequently neither invokes those results to claim a new native/cylinder bridge.

The learning obstruction is a useful \emph{precondition test} for such a bridge: an update must descend through observational equivalence before a strong operational quotient can be constructed. Its failure under GD on the flow quotient blocks that particular quotient map. It does not refute Para, the native relative operational system, or every lax, set-valued, or span-valued comparison.

\section{Comparative candidate and corrected assumptions}
\label{sec:comparative}

\subsection{Supplied data and derived structure}

Supply finite preparations $P$, controls $C$, recorded outcome sets $O_c$, terminal tests, and an exact joint response law $H(p,w,t)$ for finite event histories $w$ and terminal event $t$. Event letters are $(c,o)$; choosing a control admits \emph{all} its stochastic outcomes. Terminal tests end a run. The law includes normalization $u$. It determines adaptive public protocols by branch probabilities and grants no hidden-state interventions.

Also supply an integer rank bound $b$ and a complete exact rational Hankel block with prefix/suffix length sufficient for the standard $2b-1$ reconstruction. The global rank bound is a \emph{model-class promise}, not a conclusion of finitely many records. Classical reachable/observable linear reconstruction gives a rational row space $V$, $k=\dim V\leq b$, preparations $x_p$, event operators $A_e$, and terminal effects $\ell_t$.

Compute, rather than postulate an unknown invariant cone,
\[
 K_T=\{x\in V:\ell_t(x)\geq0\text{ for every listed terminal event }t\}.
\]
The independent finite checks are:
\begin{enumerate}
\item $K_T$ is pointed and full-dimensional; $u$ is strictly positive away from zero; the section $\Delta_T=K_T\cap\{u=1\}$ is a simplex with $k$ normalized vertices $v_i$.
\item Preparations belong to $\Delta_T$; every $A_e$ preserves $K_T$; for each control, $\sum_o A_{(c,o)}u=u$. These are vertex evaluations and linear equalities in the reconstructed space.
\item \textbf{Terminal normalization is checked on the whole simplex.} Every declared effect satisfies $0\leq\ell_t(v_i)\leq1$. For a grouped finite terminal measurement, its outcome effects sum to $u$. Equivalently, include and check complete normalized tests or each effect's complement. Bounds only on the actually reached preparations do not suffice.
\item In the computed vertex basis each event matrix $M_e$ is either positive monomial, with one strictly positive entry in every row and column, or a positive rank-one reset $M_e=b_eq_e$, with $b_e>0$ and $q_e$ a full-support probability row. The designated initial row $\alpha$ and all reset seeds have full support.
\end{enumerate}

The third check repairs an omission in the candidate working note. Positivity alone need not prevent an event value greater than one at an extrapolated vertex. If complements are instead added to the defining inequalities, the cone, its vertices, and the matrix-type checks must all be recomputed. One cannot silently add inequalities while retaining the former simplex calculation.

For an exact failure of the incomplete original checks, let $V=\RR^2$, $u(x)=x_1+x_2$, and take recorded events of two separate binary tests with $\ell_1(x)=3x_1/2$ and $\ell_2(x)=3x_2/2$. Use preparations $(2/5,3/5)$ and $(3/5,2/5)$ and the single dynamic event $I_2$. The observed table is
\[
 \begin{pmatrix}3/5&9/10\\9/10&3/5\end{pmatrix},
\]
which is valid and has rank two. Every history preserves its preparation. The displayed positivity cone is the orthant, and its normalized vertices assign probability $3/2$ to a recorded event. Thus those computed vertices are not valid stochastic states. This refutes that construction, not the existence of an alternative two-state realization. Adding complement inequalities changes the normalized interval to $x_1\in[1/3,2/3]$ and requires recomputing the checks.

The supplied interface may be diagnostically strong. Computed vertices can be unattained extrapolated laws. The output is a valid realization under the checked doctrine, not identification of the actual hidden states, a unique positive cone, or a discovered physical partition.

\subsection{Exact growth theorem}

Let $\Gamma$ be the group generated by the projective classes of the invertible monomial $M_e$. Its permutation image $G$ is finite. The kernel $D$ is a finitely generated torsion-free abelian group of positive diagonal projective matrices. Write
\[
 g=\rank_{\ZZ}D.
\]
For rational matrices, enumerate a permutation transversal, form diagonal Schreier generators, divide each diagonal by its last entry, and compute the rational rank of their prime-valuation vectors. This is effective exact algebra; factoring and permutation-group enumeration may be expensive.

Let $Q_H$ be the quotient of \emph{finite accessible histories from the designated preparation} by equality of all future experimental laws. Define $N(n)$ using histories of length \emph{at most} $n$.

\begin{theorem}[Repaired monomial and reset comparison]
\label{thm:comparative}
Under the supplied rank promise and all four finite checks above:
\begin{enumerate}
\item The simplex construction gives a $k$-state controlled positive realization of the entire preparation/test law, with valid terminal probabilities. No such finite hidden realization has fewer than $k$ states.
\item $Q_H$ is identified with the reachable normalized belief orbit in $\Delta_T$, using updates $\eta\mapsto \eta M_e/(\eta M_e\one)$.
\item $N(n)=\Theta(n^g)$.
\item $Q_H$ is finite exactly when $g=0$, equivalently when the monomial matrices have a common strictly positive projective left eigenvector. In that case its size is at most $(1+s)|G|$, where $s$ is the number of distinct reset seeds.
\item Removing whole controls, while retaining terminal separation and spanning preparations in the same response space, weakly decreases the growth degree. Merging or hiding recorded outcomes is not covered.
\end{enumerate}
\end{theorem}

\begin{proof}
The reconstructed cone and normalization give nonnegative stochastic event kernels and valid terminal effects. The native lower bound concerns the full multi-preparation/test law, whereas $N(n)$ concerns only the designated initial preparation; these are not silently identified as minima for the same initial-law problem. The Hankel rank bounds hidden cardinality below by $k$. These are standard realization facts, not additional novelty. Linear observability identifies normalized belief vectors precisely with their future laws; this is the standard filtering interpretation, not a map from a predictive state into a single hidden state.

For the upper growth bound, $D\cong\ZZ^g$ has finite index in $\Gamma$. A fixed finite transversal and Schreier generators express every word of length at most $n$ using a finite permutation choice and lattice coordinates bounded by $O(n)$. Hence there are $O(n^g)$ projective matrices. A history with resets is determined for predictive purposes by its last reset seed and subsequent monomial suffix, so finitely many seed orbits preserve this bound.

For the lower bound it is necessary to use the positive semigroup, not only the group. Let $D^+$ be the positive monomial words with identity permutation and let $H=\langle D^+\rangle\leq D$. If a monomial generator $s$ has permutation order $l$, then for $d\in D^+$,
\[
 sds^{-1}=(sds^{l-1})(s^l)^{-1}\in H.
\]
Conjugation on diagonals has finite order, so $H$ is invariant under both $s$ and $s^{-1}$. In $\Gamma/H$ every generator inverse can be replaced by a positive power, since $s^l\in H$. Any element of the diagonal kernel therefore has a positive representative with identity permutation; consequently $H=D$. Choose $g$ multiplicatively independent members of $D^+$ and fixed positive loop words representing them. Their nonnegative powers with total word length at most $n$ give $\Omega(n^g)$ distinct projective diagonals. They act freely on a full-support seed: equality of normalized images forces equality of all diagonal ratios. Observability gives distinct predictive states.

If $g=0$, the torsion-free group $D$ is trivial, so $\Gamma$ is finite. A componentwise geometric mean of a finite positive projective orbit gives a common positive projective fixed point. Conversely a common positive projective eigenvector diagonally conjugates the generators, projectively, to permutation matrices; their group is finite. Reset events add finitely many seed orbits. Restricting whole controls gives a subgroup and a subgroup of its diagonal kernel, hence a degree no larger than $g$.
\end{proof}

\paragraph{An exact parametrized family.}
For distinct primes $p_1,\ldots,p_m$, use $m$ controls, each with a monomial and a reset outcome:
\[
 M_{i,a}=\diag(1/2,1/(2p_i)),\qquad
 q=(1/2,1/2),\qquad
 M_{i,r}=\begin{pmatrix}1/2\\1-1/(2p_i)\end{pmatrix}q.
\]
Their sums are row-stochastic. From $\alpha=q$, posterior odds after a history are exactly $\prod_i p_i^{a_i}$ with $a_i\geq0$ and $\sum_i a_i\leq n$, because a reset returns odds to one. Thus
\[
 N(n)=\binom{n+m}{m}.
\]
With rank-two preparation/test access, native size is two for every $m$. Retaining $j$ whole controls gives $\binom{n+j}{j}$. This quantifies an actual distinction between two native hidden states and an arbitrarily high polynomial degree of finite-prefix predictive complexity.

Two full-support preparations and suitable terminal tests can reconstruct this example without allowing preparation of hidden pure states. One may use complementary binary tests whose event effects are $(0,1/2)$ and $(1/2,0)$. They compute the same simplex while avoiding perfect single-shot diagnosis, but they still supply strong separating exclusion information. This is not interface discovery.

A finite nontrivial control has
\[
 S=\begin{pmatrix}0&1/8\\1/2&0\end{pmatrix},\quad
 q=(1/3,2/3),\quad R=\begin{pmatrix}7/8\\1/2\end{pmatrix}q,
 \quad\alpha=(1/2,1/2).
\]
The nonreset odds map is $o\mapsto4/o$. The initial orbit is $\{1,4\}$ and the reset orbit is $\{1/2,8\}$, so the predictor has four states while the hidden model has two.

\paragraph{Boundaries.}
Full support is necessary for the lower bound: a boundary seed can be fixed by an infinite diagonal group. Separation is necessary because different coordinates can have identical future laws. ``At most $n$'' is not ``exactly $n$''; the latter can have a different degree. Finite-prefix states are not stationary infinite-past causal states. Generic positive matrices need not be monomial or resets. Exact rational multiplicative relations are unstable under perturbation, so this degree is not a robust finite-sample statistic. Hidden outcomes must be re-filtered and may leave the checked class entirely.

\subsection{Strongest reduction and Gate A status}

The project ledger already contains positive reconstruction, filtering, and finite-hidden/infinite-predictive examples. Allauzen--Mohri's twins criteria control finite weighted determinization in their stated unambiguous classes, but must not be attributed to arbitrary ambiguous automata \cite{Twins}. More strongly, Bell--Smertnig decide deterministic equivalence over number fields; the proof characterizes it by requiring every irreducible component of the minimal reachability linear-Zariski hull to have dimension at most one, meaning a finite union of lines rather than a one-dimensional ordinary linear span \cite{BellSmertnig}. Encode terminal tests by distinct terminating symbols and include normalization, then take the designated-initial reachable/observable representation. On nonzero reachable prefixes, projective equality is precisely equality of normalized future laws. The finite-union-of-lines criterion therefore settles the candidate's finite-predictor question in a broader class. This is an actual reduction, not merely related terminology.

The explicit $\Theta(n^g)$ count is not asserted by those inspected statements. It follows here from standard finite-index abelian growth plus the semigroup argument. Its absence from an inspected statement does not establish novelty. Work on arithmetic and monomial representations is a further strong comparison target \cite{PuchSmertnig}. The honest residual is a narrow quantitative application, not a new finiteness principle. Gate A's genuine-new-content requirement is not discharged.

\section{Learning candidate and corrected theorem}
\label{sec:learning}

\subsection{Experimental presentation}

Supply a labelled connected bipartite graph $H=(L\sqcup R,E)$ with $d\geq2$ vertices, a distinguished observed edge $e_0=\{a,b\}$, the Euclidean parameter metric, and the class of two-factor linear implementations $X\in\RR^{d\times w}$ with unknown finite width $w$. The architecture, layer split, row labels, metric, graph, and intervention semantics are inputs. Splitting rows into $U,V$ gives predictor $UV^{\mathsf T}$. Only $y=X_aX_b^{\mathsf T}$ is observed.

For every control edge $e=\{i,j\}$, the executable objective is the coordinate linear loss $\ell_e(X)=X_iX_j^{\mathsf T}$. Its negative Euclidean gradient is $-S_eX$, where $S_e=E_{ij}+E_{ji}$. Two distinct experimental regimes are declared:
\[
 \text{flow: }X\mapsto e^{-tS_e}X\ (t\geq0),\qquad
 \text{GD: }X\mapsto(I-\eta S_e)X\ (0<\eta<1).
\]
A single fixed positive GD step is sufficient. Independent selection and concatenation of these edge interventions are available. Sparse sensor access does not itself let one implement arbitrary squared-loss targets using unobserved predictor coordinates. The linear losses are supplied interventions, not inferred control access.

The evidence can be an ideal law on all declared protocols or the finite exact reset/replay table below. Exact resets to the same initial implementation and exact endpoint measurements are essential idealizations. No hidden-weight readout, momentum, adaptive optimizer state, nonlinear activation, regularization, passive-only training claim, or architecture change is included.

The proof variable $Z=XX^{\mathsf T}$ is not supplied as a measured state. The construction derives an observation space, actual finite probe words, the response-equivalence class, its closed updates, and compatible native implementations from the known protocol operators and observed responses. It does not discover the architecture or explain feature formation during training.

\subsection{Statement with the two-vertex exception}

Let $J=\diag(+1\text{ on }L,-1\text{ on }R)$, $D=d(d+1)/2$, and $\langle Q,Z\rangle=\tr(QZ)$. Write $\OO_{\mathrm F}$ and $\OO_{\mathrm G}$ for the spans of sensor observables pulled back through all respective finite protocols. Using $S_{e_0}$ rather than $S_{e_0}/2$ does not change the spans.

\begin{theorem}[Sparse response reconstruction and optimizer comparison]
\label{thm:learning}
For the supplied connected bipartite graph and one observed edge:
\begin{enumerate}
\item \textbf{Flow.} $\OO_{\mathrm F}=J^\perp\subset\Sym_d$. Two native initial implementations have identical responses on all flow protocols exactly when $Z'-Z=cJ$. The response state is
\[
 \QQF=\operatorname{image}\bigl(\PSD^d\longrightarrow\Sym_d/\spanlin(J)\bigr),
 \qquad q_{\mathrm F}(Z)=[Z].
\]
It has well-defined updates $[Z]\mapsto[LZL^{\mathsf T}]$, $L=e^{-tS_e}$.
\item \textbf{Sharp flow jets.} Mixed chronological response derivatives of total order at most $2d-3$ distinguish all flow-response classes. For the path $1-\cdots-d$ with sensor $\{1,2\}$, the bound is attained: $Z=0$ and $Z=E_{dd}$ agree on all lower-order mixed jets and disagree at order $2d-3$.
\item \textbf{GD.} For one fixed $0<\eta<1$, if $d\geq3$, then $\OO_{\mathrm G}=\Sym_d$ and $\QQG=\PSD^d$. Hence equality of all GD responses is precisely equality of Gram matrices. If $d=2$, then $\OO_{\mathrm G}=J^\perp$ and $\QQG=\QQF$ instead.
\item \textbf{Comparison.} For $d\geq3$, $Z\mapsto[Z]$ is a surjective information-forgetting map from GD response states to flow response states. GD updates do not descend through it in general. For $d=2$, the comparison is the identity between quotient state sets and GD updates do descend. Exact samples of the flow preserve the flow quotient in every dimension.
\item \textbf{Native width.} A flow state represented by $Z_0$ has a width-at-most-$w$ realization exactly when
\[
 \exists c\in\RR:\quad Z_0+cJ\succeq0,\qquad\rank(Z_0+cJ)\leq w.
\]
Its minimum compatible width is the minimum rank over that PSD pencil. For a GD state on a connected graph with $d\geq3$, minimum width is $\rank Z$. For $d=2$, the GD minimum is the same pencil minimum as for flow, \emph{not} the rank of an arbitrary representative.
\end{enumerate}
\end{theorem}

The last exception repairs the original working note's unconditional GD width statement. For example $I_2$ and $\diag(2,0)$ differ by $J$ and have identical GD laws with the sole edge sensor, while their ranks are two and one.

\subsection{Proof of the observable spaces}

\begin{proof}[Proof of Theorem~\ref{thm:learning}]
Set $T_e(Q)=S_eQ+QS_e$. The lifted flow is $\dot Z=-S_eZ-ZS_e$, so successive sensor Lie derivatives are words in the $T_e$, up to known signs. Endpoint responses are analytic, and equality of all nonnegative-duration protocols is equivalent to equality of these derivatives. Since $S_eJ+JS_e=0$, every generated observable is orthogonal to $J$.

For distinct $i,j,k$, let $Q_{ij}=E_{ij}+E_{ji}$. If $\{j,k\}$ is an edge, then
\[
 T_{jk}(Q_{ij})=Q_{ik}.
\]
Thus an anticommutator moves one of two distinct indices along a graph edge. The resulting two-token graph is connected when $H$ is connected. Starting at $\{a,b\}$ generates every off-diagonal $Q_{ij}$. For every edge,
\[
 T_{ij}(Q_{ij})=2(E_{ii}+E_{jj}).
\]
The edge-sum diagonal vectors span the hyperplane orthogonal to the bipartition signs: any annihilating vector satisfies $v_i=-v_j$ on each edge and is therefore a multiple of the sign vector. Off-diagonals and these diagonal sums give exactly $J^\perp$. Trace duality now identifies response equivalence with difference in $\spanlin(J)$.

For $d=2$ the jet bound is checked directly by $T_{12}(Q_{12})=2I_2$. For $d\geq3$, the two-token distance is at most $2d-4$: the known diameter bound is at most twice the graph diameter when that diameter is at most $d-2$, and the remaining diameter-$d-1$ graph is a path, whose two-token diameter is $2d-4$ \cite{Token}. One further anticommutator produces the diagonal sums. This gives the $2d-3$ jet bound. On an endpoint-sensed path, a coefficient at $E_{dd}$ requires the two indices to move from positions $1,2$ to $d,d$, requiring at least $(d-1)+(d-2)$ edge moves. Each anticommutator moves only one index one edge, including words with diagonal intermediate terms. A nonzero word of that length exists. The matrices $0$ and $E_{dd}$ are PSD and both admit width-one implementations, so the lower bound is native, not an indefinite-matrix artifact.

For GD, the observable pullback is
\[
 P_e(Q)=(I-\eta S_e)Q(I-\eta S_e)
       =Q-\eta T_e(Q)+\eta^2 R_e(Q),\qquad R_e(Q)=S_eQS_e.
\]
The eigenvalues of $S_e$ are among $1,-1,0$. On symmetric matrices, the possible eigenvalues of $P_e$, in increasing order for $0<\eta<1$, are
\[
 (1-\eta)^2,\ 1-\eta,\ 1-\eta^2,\ 1,\ 1+\eta,\ (1+\eta)^2.
\]
They are distinct; on their eigenspaces $T_e$ has eigenvalues $2,1,0,0,-1,-2$, respectively. Spectral interpolation makes $T_e$ a polynomial in $P_e$. Hence the GD cyclic space contains the flow cyclic space and is invariant under $R_e=(P_e-I+\eta T_e)/\eta^2$. Missing eigenspaces when $d=2$ do not affect the argument.

If $d\geq3$, connectedness supplies adjacent edges $ij,jk$. The flow space contains $E_{ii}+E_{jj}$, and applying $R_{jk}$ produces $E_{kk}$, which is not in $J^\perp$. Thus GD generates all of $\Sym_d$. If $d=2$, instead
\[
 (I-\eta S_e)J(I-\eta S_e)=(1-\eta^2)J,
\]
so the invisible line is preserved and the GD space remains exactly $J^\perp$.

For flow, $LJL^{\mathsf T}=J$ proves update descent and the closed quotient dynamics. For GD with $d\geq3$, a counterexample is given below. A PSD matrix has a factorization $XX^{\mathsf T}$ with $w$ columns exactly when its rank is at most $w$. This proves all width claims, with the quotient rather than an arbitrary Gram representative used in the two-vertex case.
\end{proof}

\subsection{Typed comparison and an exact failure of descent}

For $d\geq3$ the comparison factors through common native implementations:
\[
 \coprod_{w\geq0}\RR^{d\times w}
 \xrightarrow{\ X\mapsto XX^{\mathsf T}\ }\QQG=\PSD^d
 \xrightarrow{\ q_{\mathrm F}\ }\QQF.
\]
This is an information map between laws of two different dynamics. GD protocols are not a subset of flow protocols, and polynomial interpolation of observables is not a physical simulation of one protocol by another. The map does not claim an operational morphism intertwining both dynamics.

For the path $1-2-3$, take $J=\diag(1,-1,1)$, sensor $y=Z_{12}$, and $Z_0=I_3$, $Z_1=I_3+cJ$ with $0<|c|<1$. Both are PSD. Their flow laws are identical. After a GD step on edge $23$ their difference is
\[
 c\diag(1,-1+\eta^2,1-\eta^2),
\]
which is not a multiple of $J$. Thus that update already separates their flow-response classes. A subsequent GD step on edge $12$ makes the fixed sensor difference exactly $-c\eta^3$. Consequently the dashed bottom arrow in
\[
\begin{tikzcd}[column sep=large]
\PSD^3 \arrow[r,"U^{\mathrm G}_{23}"] \arrow[d,"q_{\mathrm F}"'] &
\PSD^3 \arrow[d,"q_{\mathrm F}"]\\
\QQF \arrow[r,dashed,"\text{no such }\overline U^{\mathrm G}_{23}"'] & \QQF
\end{tikzcd}
\]
cannot make the square commute. No claim is made that all other comparison classes are impossible.

At $c=-1$, $I_3$ and $\diag(0,2,0)$ still have the same flow law, but their GD minimum widths are three and one. The common flow class has minimum compatible width one. This is information loss with a native-cost consequence, not physical deletion of columns from a width-three implementation.

The width contrast also has an unbounded family. On a star with $d\geq3$, put $J=\diag(1,-1,\ldots,-1)$. The Grams $I_d$ and $I_d+J=\diag(2,0,\ldots,0)$ have the same nonzero flow law, hence common minimum compatible flow width one, but their GD minimum widths are $d$ and one. This is an observation/resource contrast across a graph family, not a growing representation during one training run.

\subsection{Native fibres and disconnected systems}

Actual implementation width is not identifiable even in the full-Gram regime: appending zero columns changes width without changing any response. All minima here are response-compatible native widths.

The feasible parameters $c$ of a nonempty PSD pencil form a compact interval: positive diagonal entries bound $c$ in both directions because both bipartition classes are nonempty. If the interval has distinct endpoints, every strict interior point is positive definite. Indeed a kernel vector of an interior PSD combination is in both endpoint kernels and hence is killed by their invertible difference $(c_2-c_1)J$. Minimum rank is therefore attained at an endpoint. A singleton interval requires its own rank calculation. At fixed width, factorizations of a fixed Gram differ by a right orthogonal transformation after zero-padding; different pencil points are additional behaviorally invisible choices, not that same gauge.

For a disconnected graph with a sensor in every nontrivial component, the response laws identify the component diagonal blocks only. Cross-component Gram blocks remain invisible. Flow identifies each block modulo $\RR J_C$. GD identifies a block exactly when $|C|\geq3$, and modulo $\RR J_C$ when $|C|=2$. Isolated vertices are unobserved and can be set to zero in a minimum implementation.

The minimum shared hidden width is the \emph{maximum} of the component minima, not their sum. For the upper bound, factor each block in the same sufficiently large $\RR^w$ and stack the factors; their unobserved cross-block products are unconstrained. The reverse bound follows by restricting a global Gram factorization to each block. Independent feature pools, prescribed cross-block correlations, or extra joint tests would define a different native problem.

\section{Finite experiments and what their cost means}
\label{sec:experiments}

\begin{proposition}[Exact finite endpoint reconstruction]
\label{prop:endpoints}
From the graph and sensor alone one can construct a finite list of endpoint protocols, independent of the unknown weights and width, whose exact reset responses reconstruct the response state in Theorem~\ref{thm:learning}. For connected $H$:
\begin{enumerate}
\item flow needs $D-1$ scalar responses, using only edge durations $\tau=\log2$, with at most $D-2$ segments per selected protocol;
\item GD with one fixed rational $0<\eta<1$ needs $D$ scalar responses and at most $D-1$ steps per protocol when $d\geq3$;
\item for $d=2$, both need $D-1$ scalar responses, identifying the quotient rather than the Gram.
\end{enumerate}
These are sufficient finite counts in the supplied model class, not sample-complexity or globally optimal experimental-cost claims.
\end{proposition}

\begin{proof}
At $\tau=\log2$,
\[
 L_e=e^{-\tau S_e}=I+\tfrac14S_e^2-\tfrac34S_e
\]
has rational entries. Its observable pullback is $E_e(Q)=L_eQL_e=e^{-\tau T_e}Q$. The eigenvalues of $T_e$ belong to $\{-2,-1,0,1,2\}$; exponentiation at this fixed nonzero duration separates them, so $T_e$ is a polynomial of $E_e$ of degree at most four. Thus the fixed-duration cyclic space equals the full flow space. The GD interpolation in the previous proof gives its corresponding equality.

Start with the sensor matrix and repeatedly apply each known pullback to the current basis words. Keep a word only when its rational matrix vector increases rank. When all basis images lie in the span, the span is invariant and complete. If its final dimension is $m$, no retained basis word needs length more than $m-1$. The retained matrices give $m$ linearly independent trace measurements. Run those actual words with a reset before each run and solve their known linear system. The dimensions are $m=D-1$ for flow and for the two-vertex GD case, and $m=D$ for GD when $d\geq3$.
\end{proof}

The word acts by chronological state updates; its observable matrix is computed by pullback in the corresponding reversed order. The scalar readout is $\frac12\tr(S_{e_0}Z)$, so the known factor of two must be included in the measurement matrix. A breadth-first implementation gives the stated word-length bounds without assuming derivatives are directly observable.

A simple closure algorithm tests $O(|E|D)$ generated candidates, each requiring known matrix arithmetic and a rank test. This is a coarse exact-arithmetic construction, not an optimized bit bound. The $D-1$ or $D$ observations each need a fresh exact reset, an exact scalar readout, and the specified sequence of controlled training operations. Total execution length is at most quadratic in $D$ under these loose bounds; fixed flow duration also has a nonzero physical cost. Measurement precision, preparation reproducibility, and implementation of linear losses are separate resources.

\paragraph{Model falsification versus certification.}
A GD table reconstructs a candidate Gram and can be rejected if it is not PSD. A flow table reconstructs a class and can be rejected if its PSD pencil is empty. Held-out protocols can test the stipulated class. Matching the finite basis table certifies the entire law only \emph{conditional on the known model class and update law}; arbitrary systems can match that table and disagree later.

Noise produces an uncertainty set under a known linear map. It does not identify exact rank or a pencil endpoint without conditioning and margin assumptions. The displayed GD discriminator has magnitude $|c|\eta^3$. Identifiability for every $\eta>0$ is not uniform stability as $\eta\downarrow0$. No finite-sample guarantee, confidence coverage, physical duration optimization, or robust native-width estimator is proved.

\paragraph{Comparative experimental cost.}
The positive-realization candidate starts from a complete exact rational finite response block plus a global rank promise. Its full branching block can grow exponentially with the horizon and alphabet. Cone computations, rational factorization, and finite permutation enumeration are effective but not asserted polynomial-time. Exact support, exact rank, and exact multiplicative dependence are not finite-noisy-data observations. Native hidden count, predictive state count, arithmetic description size, probe count, and sampling cost are distinct resources.

\section{Strongest prior art and the residual actually left}
\label{sec:sources}

\subsection{Learning reconstruction is classical before the graph calculation}

Sontag's observable-space construction and canonical realization results already generate state information from controlled future responses and output jets. Hermann--Krener supply the differential observability framework. Grossman--Larson's bialgebraic input-output realization controls formal-series recodings. A response quotient or the smallest update-invariant observable algebra is therefore not a new learning principle \cite{Sontag,HermannKrener,GrossmanLarson}.

After lifting to $Z$, this candidate is a finite bilinear controlled system. Petreczky's Theorem 2.2(ii) and Remark 2.1 express indistinguishability by kernels of output matrix words and give observation/minimization constructions; Theorem 2.3 characterizes minimal realizations by reachability, observability, and Hankel rank \cite{Petreczky}. Proposition~\ref{prop:endpoints} is a direct cyclic-module application once the graph space is known. Calling it a finite Koopman lift does not add a theorem; finite invariant observable spaces have their own established theory and limitations \cite{Koopman}.

For nonlinear neural predictors, the first learning response depends on a parameter-Jacobian Gram operator, as in the NTK construction. The neural tangent hierarchy cautions against assuming that adjoining one such matrix closes general finite-width nonlinear dynamics \cite{NTK,NTH}. The exact degree-two closure here is special to the supplied linear architecture and losses.

\subsection{The direct matrix-learning predecessor}

Zhang--Peyr\'e is a particularly strong, recent control, not an optional distant citation \cite{ZhangPeyre}. Proposition 2.1 supplies exact covariance closure. Definition 2.2, equation (5), and Proposition 2.3 supply the signed PSD interval, spectral endpoints, and invariance. Appendix D, Proposition D.1, gives generation of the indefinite-orthogonal Lie algebra by constant-gradient split fields; Corollary D.2 describes generic common smooth conservation laws locally through the signed spectrum. Lemma C.14 and equation (42) explicitly treat the Euler signature defect.

Consequently none of the following is a residual: introducing the covariance $XX^{\mathsf T}$; preserving $J$ under flow; losing that exact preservation under Euler steps; constructing the PSD pencil; using its endpoint ranks; or observing that full-control Gram dynamics close. The graph-restricted, one-sensor module calculation, its sharp path jet bound, and the exact information comparison are the remaining statements written out here. The inspected text does not explicitly state this combined theorem. That observation is \emph{not} evidence that it is absent as a corollary, known elsewhere, or publication-worthy. The full Lie-algebra and bilinear-observability baselines make an elementary reduction plausible.

The token-graph connectivity and diameter estimates used in the proof are prior art \cite{Token}. The direct Euler signature-defect source already prevents claiming the optimizer distinction from conservation alone as new. The candidate notes also identify a discretization/observability paper in an unmeasured-input setting \cite{Discretization}; its primary text could not be independently reverified in the concluding audit, so no precise result from it is needed for this decision.

\subsection{What may still be a useful theorem without being new}

The corrected learning result makes a clean operational distinction: a finite family of sparse interventions reveals every Gram observable except one direction under flow, and all of them under GD on connected graphs of size at least three. It tells an experimenter which state is sufficient, gives actual finite probes, identifies a sharp response order, and exposes a failed quotient update. The finite native family and width loss are concrete. These are worth retaining as a benchmark and a candidate instance of the accepted learning target.

However, novelty cannot be manufactured by grouping classical consequences under one title. A real Gate A pass would need a source-supported argument that a substantive premise/conclusion obligation remains beyond those predecessors, or a sharpened problem whose solution requires additional mathematics. At present the report establishes the exact obligation to inspect, not that it is new.

\section{Noncircularity and falsifying controls}
\label{sec:controls}

\subsection{Supplied versus derived}

The comparative theorem tests a \emph{computed} cone, vertices, terminal effects, and matrices. It does not assume the existence of an arbitrary sought invariant cone or finite predictor. Its exact-law/rank and strong diagnostic promises are nevertheless substantial supplied structure. The positive realization is not empirically forced among every possible native doctrine.

The learning theorem assumes a graph, a two-factor architecture, known update rules, linear-loss actuation, a scalar sensor, and reset access. It does not assume the resulting observable-space dimension or provide an unknown sufficient state as input. Connectedness plus the explicit matrix laws proves that dimension. The unknown Gram or pencil is recovered from responses. This is a noncircular construction \emph{within the supplied class}, but the hidden features, architecture, and metric have not emerged from unrestricted observations.

\subsection{Controls required by the claimed scope}

\begin{description}[style=nextline,leftmargin=0pt]
\item[Passive training and stationary loss]
One passive equilibrium need reveal no extra Gram information. Independent linear-loss interventions are essential. Neither theorem applies to a claim based only on loss curves or observationally unavailable controls.
\item[Coordinate and optimizer changes]
A right orthogonal hidden-feature transformation preserves $Z$ and the declared Euclidean dynamics. A general invertible feature rescaling can preserve current predictors but changes ordinary Euclidean learning unless the metric and vector fields are transported too. GD is a different learner from exact sampled flow. Hidden momentum or adaptive state invalidates a Gram-only state.
\item[Architecture and formation]
Neither theorem discovers the graph, partition, feature grammar, or decomposition. In the invertible update regime, native Gram rank is constant. Flow also preserves the feasible pencil ranks under congruence. Response reconstruction is not representation formation or a proof that new physical parts appear during training.
\item[Sensor and support boundaries]
Unsensed connected components and cross-component Gram blocks are invisible. For the stochastic candidate, boundary seeds and insufficient terminal separation can collapse predictive distinctions. The terminal normalization repair is required on extrapolated simplex vertices, not just observed states.
\item[Probability versus possibility]
Support incidence can stay fixed while arithmetic predictive growth changes. O7-style possibility data and quantitative law data remain distinct. No FCA completion recovers omitted numerical information merely by naming it as a concept.
\item[Finite evidence]
Both finite reconstruction procedures are conditional on their model-class promises. Noise, finite precision, inaccurate resets, and an untested continuation require uncertainty or stability analysis. The proofs do not provide it.
\item[Native and joint semantics]
The learning width formula is a Gram factorization criterion; the comparative hidden cardinality bound is a Hankel-rank criterion. Neither is a retained-mechanism refinement theorem, a same-budget lift for every coarse native model, or independent multistage assembly from partial local witnesses. Unobserved cross-block correlations must not be inferred from marginal reconstruction.
\item[Comparison direction]
Positive realization maps to a space of belief distributions before predictive minimization; there is no canonical map from each predictive state to a single hidden state. GD-to-flow is an information map on common native Grams, not a protocol inclusion. A failed update descent does not eliminate valid one-sided or weaker comparisons.
\item[Feedback and approximate claims]
The stochastic law determines adaptive branch probabilities, but the growth theorem covers finite histories in its restricted matrix class. No feedback stability bound or approximate compositional theorem follows. Learning is treated as deterministic controlled words, not a statistically sampled adaptive acquisition policy.
\end{description}

Mandatory project controls concerning independent joint assembly, retained stochastic anchors, native optimization, and noisy feedback remain outside these theorem statements. They are marked open rather than declared passed vacuously. No whole-programme generality is claimed.

\section{Verification status and the next honest decision}
\label{sec:status}

The report separates four levels:
\begin{enumerate}
\item \textbf{Mathematical proof seeds.} Theorems~\ref{thm:comparative} and~\ref{thm:learning}, with the terminal-effect and two-vertex repairs, have direct proofs in this report. The endpoint construction and width consequences are justified at their stated exact scope.
\item \textbf{Exact finite checks.} The accompanying rational scripts check finite graph instances, the $-c\eta^3$ discriminator, selected predictive counts, and boundary controls. These detect errors and reproduce examples; they are not proofs of the universal theorems or formal verification.
\item \textbf{Independent review and source comparison.} The accompanying reviews found the two substantive statement defects now repaired. Source inspection supplies exact prior-art reductions and distinguishes inspected statements from broad literature claims. Independent review is not novelty clearance.
\item \textbf{Unproved research claims.} A new central Gate A residual, a general cross-mechanism bridge, representation formation, empirical acquisition guarantees, native multistage assembly, and a relative/cylinder reflection theorem remain unestablished.
\end{enumerate}

The original challenge script checks path jet dimensions for $2\leq d\leq8$, endpoint-flow and one-/two-step-size GD dimensions for paths through $d=6$, and selected stars through $d=6$. It distinguishes jet depth from endpoint word depth. Its finite bounds do not prove the all-graph theorem. The independent standard-library rational bundle checks 54 graph-closure cases, five fixed-step interpolation identities on all symmetric $3\times3$ basis matrices, star width contrasts for $3\leq d\leq12$, the two-vertex exception, and the terminal-normalization counterexample and repair. It enumerates prime-control predictive orbits for $m=1,2,3,4$ through $n=8$, obtaining respectively $9,45,165,495$ states at $n=8$, and checks the four-state finite orbit and exact-time/at-most distinction. Normal and optimized Python runs produce byte-identical results; deliberate negative controls fail in both. The bundle uses no floating-point arithmetic or assertion statements. Its manifest defines the precise coverage and hashes.

\paragraph{Gate A assessment.}
Precise semantics, independently checkable class conditions, positive examples, sharp failures, and effective exact reconstruction are present. The comparative candidate fixes genuinely unlike mechanisms but does not establish a new sufficient residual. The learning candidate fixes materially different dynamics, which the learning target allows, but not unlike reconstruction mechanisms in the broader comparative sense. Its strongest surviving calculation is explicit, and its novelty is unresolved with serious reasons to expect a short classical reduction. Therefore neither is promoted to a new central result.

\paragraph{Stop and repair conditions.}
Stop advertising either claim as new if the full strongest-source reduction supplies it. Stop the stated mathematical theorem if a verified counterexample survives all corrected hypotheses. If the only repair inserts the unknown observation space, invariant cone, or desired native lift as an assumption, reject it as circular. If required interventions or resets are unavailable, stop the empirical claim rather than silently replacing them with an oracle. A narrower theorem or benchmark can remain useful without fulfilling a stronger programme commitment.

\paragraph{Programme stance.}
The outcome favors a disciplined family of mechanisms with explicit bridges over an unsupported universal wrapper. It does not decide that this must be the final architecture of the theory. The accepted ambition still needs a genuinely added result relating observed distinctions, generated relations, predictive continuation, and native admissibility in their correct types. The learning example can test a response-sensitive branch; it cannot stand in for that whole question. No additional project stage, publication claim, recurring work, or alteration of the accepted contract is authorized by this assessment.

\appendix
\section{Artifact and source discipline}

The main source is \texttt{report.tex}; the README states the proof status, execution commands, correction hierarchy, and absence of compilation. The candidate working notes are provenance, not the final corrected theorem statements. The accepted contract source has been left unchanged. No upload, PDF compilation, or rendered-layout check is claimed.

References to earlier programme results describe their status in the accepted contract and recovered ledger, not fresh re-verification. Exact theorem attributions below are controlled by the accompanying source audit. Abstract-level or phenomenon-level comparisons are not promoted to a full literature exclusion. Source search cannot prove absence of prior art.

\begin{thebibliography}{99}
\bibitem{Contract} \emph{A research contract for categorical phenomenology}, revised discussion draft with parametric amendment, 3 November 2026. Accepted project baseline supplied for this assessment; unchanged here.
\bibitem{Ledger} \emph{Recovered phenomenology results: theorem ledger}, 2 November 2026. Programme source recovery, covering generated images, Boolean/stochastic comparisons, prediction, positive Hankel realization, and completed cylinder/cross-doctrine checkpoints.
\bibitem{Relative} \emph{One relative construction for context content and effects}, 1 November 2026, and \emph{Integrated Relative Computation Acceptance Dossier}, 2 November 2026. Programme dependencies for O1, O5, and O7; their original hypothesis and review status is retained.
\bibitem{FCA} R. Wille, \emph{Restructuring Lattice Theory: An Approach Based on Hierarchies of Concepts}, 1982. \url{https://doi.org/10.1007/978-94-009-7798-3_15}. Concept generation is a prior-art control; the project ledger records its specific row/meet comparison.
\bibitem{Stacey} A. Stacey, \emph{Comparative Smootheology}. \url{https://arxiv.org/abs/0802.2225}. Definitions 3.1, 3.4, 3.6, 3.8 and Corollary 4.16, with their adequacy/forcing assumptions; stronger comparisons are separately assessed in the source.
\bibitem{Garner} R. Garner, \emph{The Isbell monad}. \url{https://arxiv.org/abs/1410.7108}. Theorems 3.3, 4.1, and 5.1.
\bibitem{CALF} G. van Heerdt, M. Sammartino, A. Silva, \emph{CALF: Categorical Automata Learning Framework}. \url{https://drops.dagstuhl.de/entities/document/10.4230/LIPIcs.CSL.2017.29}. Theorem 9, within its wrapper and factorization hypotheses.
\bibitem{Twins} C. Allauzen, M. Mohri, \emph{Efficient Algorithms for Testing the Twins Property}. \url{https://cs.nyu.edu/~mohri/pub/twins.pdf}. Proposition 1, Theorems 5--6, and the real/log transfer with its unambiguity restrictions.
\bibitem{BellSmertnig} J. P. Bell, D. Smertnig, \emph{Computing the linear hull: Deciding Deterministic? and Unambiguous? for weighted automata over fields}, LICS 2023. \url{https://math.smertnig.at/paper/linear_hull.pdf}. Theorem 1 and its reachability-linear-hull proof.
\bibitem{PuchSmertnig} A. Puch, D. Smertnig, \emph{Factoring through monomial representations: arithmetic characterizations and ambiguity of weighted automata}. \url{https://arxiv.org/abs/2410.03444}. Further arithmetic/monomial reduction target, not an identification of path ambiguity with predictive-state count.
\bibitem{Sontag} E. D. Sontag, \emph{Spaces of Observables in Nonlinear Control}. \url{https://sontaglab.org/FTPDIR/icm.pdf}.
\bibitem{HermannKrener} R. Hermann, A. J. Krener, \emph{Nonlinear Controllability and Observability}. \url{https://www.math.ucdavis.edu/~krener/1-25/10.IEEETAC77.pdf}. Differential-observability baseline; not attributed the exact global graph theorem here.
\bibitem{GrossmanLarson} R. Grossman, R. G. Larson, \emph{The realization of input-output maps using bialgebras}. \url{https://arxiv.org/abs/2007.09526}. Formal realization control, not an empirical reconstruction guarantee.
\bibitem{Petreczky} M. Petreczky, \emph{Realization theory for linear and bilinear switched systems: A formal power series approach}, ESAIM COCV, 2011. \url{https://www.numdam.org/item/10.1051/cocv/2010015.pdf}. Theorems 2.2 and 2.3, Remark 2.1.
\bibitem{Koopman} S. L. Brunton and collaborators, \emph{Koopman Invariant Subspaces and Finite Linear Representations of Nonlinear Dynamical Systems for Control}. \url{https://arxiv.org/abs/1510.03007}. Used as a finite-observable-space control, not a theorem-specific attribution for this graph class.
\bibitem{NTK} A. Jacot, F. Gabriel, C. Hongler, \emph{Neural Tangent Kernel: Convergence and Generalization in Neural Networks}. \url{https://arxiv.org/abs/1806.07572}.
\bibitem{NTH} J. Huang, H.-T. Yau, \emph{Dynamics of Deep Neural Networks and Neural Tangent Hierarchy}. \url{https://arxiv.org/abs/1909.08156}.
\bibitem{ZhangPeyre} S. Y. Zhang, G. Peyr\'e, \emph{Global Exponential Convergence of Two-Layer Linear Network Training}, arXiv:2610.09356v1. \url{https://arxiv.org/html/2610.09356v1}. Proposition 2.1, Definition 2.2, equation (5), Proposition 2.3, Lemma C.14/equation (42), and Appendix D.
\bibitem{Token} R. Fabila-Monroy and collaborators, \emph{Token Graphs}. \url{https://arxiv.org/abs/0910.4774}. Connectivity, diameter bounds, and the path calculation used in the jet proof.
\bibitem{Discretization} M. N. Chatzis, K. Maes, G. Lombaert, \emph{Observability of nonlinear discrete-time systems with unmeasured inputs}, 2025. \url{https://doi.org/10.1016/j.ymssp.2025.113142}. A phenomenon-level discretization comparison with different input hypotheses.
\bibitem{GradientLearning} G. S. H. Cruttwell, B. Gavranovi\'{c}, N. Ghani, P. Wilson, F. Zanasi, \emph{Categorical Foundations of Gradient-Based Learning}. \url{https://arxiv.org/abs/2103.01931v2}.
\bibitem{Lenses} G. S. H. Cruttwell, B. Gavranovi\'{c}, N. Ghani, P. Wilson, F. Zanasi, \emph{Deep Learning with Parametric Lenses}. \url{https://arxiv.org/html/2404.00408v1}. Includes explicit stateful optimizer constructions.
\bibitem{CDL} B. Gavranovi\'{c} and collaborators, \emph{Position: Categorical Deep Learning is an Algebraic Theory of All Architectures}. \url{https://arxiv.org/html/2402.15332v2}. Definition G.1, Theorem G.10, and Conjecture G.12 are kept distinct.
\end{thebibliography}
\end{document}
